\section{Limitations}
\label{sec:limitations}

We state the boundaries of this work explicitly, both because they define what
the reported numbers may be used for and because several of them are the natural
next steps.

\paragraph{L1. No physical validation, and an unquantified simulation-to-reality
gap} Every result in this paper was produced in simulation. No experiment was
run on a physical BlueROV2 or in a water tank, and we make no claim that any
reported completion rate, time or collision count would be reproduced by a real
vehicle. The gap is unquantified rather than small: closing it would require a
physical gate course, a calibrated acoustic positioning setup and a
characterization of the optical conditions, none of which is attempted here.

\paragraph{L2. Backend fidelity bounds the physical claims} The benchmark inherits
the hydrodynamics, rendering and sensor models of its backend. HoloOcean provides
a usable and widely adopted approximation, but it is not a validated
computational-fluid-dynamics model of a BlueROV2, and the current field of
\eqref{eq:current_sum} is an analytic superposition of primitives rather than a
simulated flow. Conclusions about \emph{relative} controller behaviour are
therefore better supported than conclusions about absolute achievable speeds or
about the specific current magnitude at which a controller fails.

\paragraph{L3. Optical perception assumptions} The detector of
Section~\ref{sec:perception_detect} assumes gates that project to closed,
near-rectangular contours of bounded aspect ratio against a background whose
edge structure is less rectangular. It was developed for, and validated on, the
rendering conditions of these circuits. Turbidity, marine snow, backscatter from
an active light source, biofouling on the gate frames, strong caustics and
low-light operation are not modelled, and the detector's geometric filters would
require re-derivation under any of them.

\paragraph{L4. Gate appearance assumptions} All gates in this study are square
frames of identical $1.5\,\text{m}\times1.5\,\text{m}$ inner aperture, and the
detector's size, aspect and near-field area gates are calibrated to that
geometry. The known-aperture assumption is also what makes the pose extension's
known-size distance estimate possible. Heterogeneous gate sizes, non-square
apertures, partially occluded gates or gates that must be distinguished by
appearance rather than by acoustic association are outside the current scope.

\paragraph{L5. Acoustic sensing is an approximation, not a channel model} The
beacon model of \eqref{eq:beacon} applies independent Gaussian perturbations to
range, bearing and elevation with Bernoulli dropout and a range cut-off. It does
not model multipath, ray bending, Doppler on the acoustic carrier,
frequency-dependent absorption or the correlated errors these produce. The
inter-vehicle channel of \eqref{eq:comms_latency} is likewise an
acoustic-inspired transport with range-dependent latency and independent packet
loss; its loss probability, range, latency and freshness window were not
independently swept, and no claim about acoustic communication performance is
made.

\paragraph{L6. The current sweep is narrow} Quantitative current results are
restricted to two named profiles on one circuit (Horseshoe Bay). The
\texttt{strong} profile turned out to be a saturation condition rather than a
graded difficulty, so the informative disturbance evidence rests effectively on a
single profile and five seeds. A finer sweep, and the same sweep on the other two
circuits, would be needed to characterize the degradation curve rather than two
points on it.

\paragraph{L7. Two baseline controller families} The rule-based evidence covers
exactly two controllers that differ in one design decision. This is deliberate ---
it makes the comparison a clean ablation --- but it means the benchmark has not
yet been exercised against structurally different autonomy strategies, such as a
model-predictive controller, a cascaded design with an inner DVL velocity loop,
or a planner that reasons about several gates ahead. We therefore cannot claim
that the difficulty ordering across the three circuits is a property of the
circuits rather than of these two controllers.

\paragraph{L8. Limited fleet scale and diversity} Fleet evidence covers two
homogeneous vehicles at a $90$\,s gap and three heterogeneous vehicles at two
gaps, on one clean circuit, with three or five seeds. Larger fleets, fleets under
current or with obstacles, formation-keeping objectives, cooperative overtaking,
dynamic role assignment and competitive multi-team racing are all unaddressed.
The coordination result in particular rests on three seeds per condition.

\paragraph{L9. The pose-aware perception extension is not validated} As reported
in Section~\ref{sec:perception_pose}, the exploratory pose front end recovers a
pose in only $11.7\%$ of frames on a dedicated oblique-view rig, with a median
plane-yaw error of $48.2^{\circ}$ and orientation-class and sign accuracies of
$0.00$. It is accurate in the near-frontal regime the controller actually
operates in (Table~\ref{tab:perception_audit}) and unreliable outside it. It is
implemented and released, but it supports no claim in this paper and should not
be used as a pose estimator.

\paragraph{L10. The learned-control results are a snapshot of an open problem}
The learned policies reported here are feed-forward PPO and behaviour-cloning
checkpoints trained on procedural geometry. The frozen policy failed a
pre-registered readiness gate on four of nineteen criteria and completed $0$ of
$30$ official-circuit holdout trials, and the best checkpoint of the matched
benchmark reaches $40\%$ on those circuits. We report a diagnosis --- the deficit
is in ordered-sequence management rather than in gate-following --- and not a
solution. Architectures with explicit memory, imitation-based bootstrapping and
interactive dataset aggregation are natural responses to that diagnosis and are
not evaluated here. A separate line of work along those lines was in progress
while this manuscript was written and is deliberately excluded: no result from it
appears anywhere in this paper.

\paragraph{L11. The official circuits are no longer unseen} An exploratory
holdout evaluation opened Horseshoe Bay, Vertical Serpent and Mixed Endurance,
and its access ledger records all sixty accesses. They must therefore not be
treated as held-out circuits for any policy trained after that point, and every
learned-controller result on them in this paper is labelled a retrospective
diagnostic benchmark. Future generalization claims require a sealed circuit set
that has not been opened.

\paragraph{L12. Three circuits do not span underwater autonomy} The official
track set was designed to vary length, verticality and sequencing, and it does so
--- 12 to 22 gates, $93.8$ to $206.3$\,m, planar to serpentine. It nonetheless
covers a narrow slice of the underwater autonomy problem. Docking, manipulation,
inspection, long-horizon navigation with drift, adversarial or dynamic
environments and mission-level planning are all outside the benchmark's task
definition, and results here should not be read as evidence about them.
